\section{Inference: derivations and implementation notes}
\label{app:inference}

This appendix supplies the derivations and implementation detail behind
\S\ref{sec:belief}. Notation is that of \S\ref{sec:notation}. The source files
are \filepath{engine/src/belief.hpp} (constraint bookkeeping and the Fast path),
\filepath{engine/src/blockdp.hpp} (the exact reference engine), and
\filepath{engine/src/oracle.hpp} (the brute-force enumerator used in E15).

\subsection{The layer identity}
\label{app:layer}

Order the unresolved cards $c_1,\dots,c_{|U|}$ and let $q = (q_0,\dots,q_5)$ be
the capacity vector of \eqref{eq:capacity}. A partial assignment of the first $j$
cards leaves a vector $n$ of unused capacities. Because every one of the first
$j$ cards consumes exactly one unit of capacity from exactly one seat,
\[
  \sum_p n_p \;=\; \sum_p q_p - j \;=\; |U| - j ,
\]
so $j$ is recoverable from $n$ and the layer index need not be stored. This is
the identity that collapses the natural $|U|$-layer table into two single arrays.
Concretely, the forward and backward tables are laid out as one array each of
size $\prod_p (q_p+1)$ --- at most $10^5$ entries by \eqref{eq:bound}, so the
$|U|$-layer table that \eqref{eq:fwd}--\eqref{eq:bwd} would otherwise require
never has to be materialised --- addressed by the mixed-radix offset
$\mathrm{off}(n) = \sum_p n_p \, s_p$ with strides $s_p = \prod_{r<p}(q_r+1)$,
and a companion index sorts the reachable states by $\sum_p n_p$ so that a sweep
visits one layer at a time. Count vectors are packed one byte per seat into a
single \verb|uint64_t|, which lets the dominance test $n \ge t$ that guards a
block transition be evaluated as one branch-free word operation.

The recursions are those of \eqref{eq:fwd}--\eqref{eq:bwd}, evaluated with the
layer constraint $\sum_p n_p = |U| - j$ in force at every step and initialised at
$F_0(q) = 1$ and $B_{|U|}(\mathbf{0}) = 1$; $e_p$ is the $p$-th unit vector and
$q$ the full capacity vector of \eqref{eq:capacity}, so a transition consumes one
unit of one seat's capacity and the branching factor at card $c_j$ is
$|M_{c_j}| \le 5$. The per-card marginals of
\S\ref{sec:counting} are the contraction of the two tables at card $j$,
\begin{equation}
  \mu_{c_j, p}
    \;=\; \frac{\mathbf{1}[\,p \in M_{c_j}\,]}{Z}
      \sum_{\substack{n \,:\, |n| = |U| - j + 1 \\ n_p \ge 1}}
        F_{j-1}(n)\; B_j(n - e_p),
  \label{eq:marg}
\end{equation}
which is \emph{exact}, since every consistent assignment is counted once by
splitting it at card $j$ into a prefix counted by $F_{j-1}$ and a suffix counted
by $B_j$. Both tables are evaluated in double
precision on integer-valued data. Each entry is a count of assignments, so the
arithmetic is exact while the counts stay below $2^{53}$, and the state-space
bound of \S\ref{app:bound} keeps the number of terms per entry at six. E15 does
not report the magnitudes of the partition functions it compared; what it
reports is that the block engine and the brute-force enumerator agreed exactly,
with a maximum relative difference of zero in $Z$ and a maximum absolute
difference of zero across all \numOracleMargChecks{} marginal,
\numOracleTeamChecks{} team-ownership and \numOracleAllocChecks{} allocation
comparisons (\S\ref{sec:results-verify}), that is, at integer rather than
rounding scale.

\subsection{Why the state space is bounded by $10^5$}
\label{app:bound}

The observer's own cards are resolved by (C1), so $q_i = 0$ and the table is a
product over the five other seats. Also $|U| \le 45$: at the deal the 54 cards
lie in nine live half-suits, the observer holds nine of them and those are
resolved by (C1), so $|U| = 45$ initially; thereafter a card leaves $U$ when its
holder is revealed or its half-suit is declared and, by
Proposition~\ref{prop:transfer}, never re-enters. Subject to $\sum_{p\ne i} q_p =
|U|$, the product $\prod_{p \ne i}(q_p+1)$ is maximised when the capacities are
equal, by the arithmetic--geometric mean inequality, giving
\begin{equation}
  \prod_{p\ne i}(q_p+1) \;\le\; \Bigl(\frac{\sum_{p \ne i}(q_p+1)}{5}\Bigr)^{5}
  \;=\; \Bigl(\frac{|U|}{5}+1\Bigr)^{5} \;\le\; 10^{5}.
  \label{eq:bound}
\end{equation}
Each state has at most six outgoing transitions, so a
full forward--backward pass is bounded by $6 \times 10^5$ multiply--adds
independently of the position, and the two arrays together occupy under two
megabytes. In practice most positions are far smaller, because resolved cards
leave $U$ and completed half-suits leave $\mathcal{H}$. The implementation
additionally refuses to build a table above a fixed state cap and reports the
failure to the caller rather than silently degrading; the build-failure count was
zero in both E2 and E15.

\subsection{The direct sampler}
\label{app:sampler}

The backward table doubles as a sampler over assignments. Having placed
$c_1,\dots,c_{j-1}$ and arrived at the remaining-capacity vector $n$, card $c_j$
is given to seat $p$ with probability proportional to $B_j(n - e_p)$, restricted
to $p \in M_{c_j}$ with $n_p \ge 1$. A seat receives positive weight only when
the suffix can still be completed from $n - e_p$, so every prefix the sampler
constructs is extendable and no draw is ever discarded: the procedure is
rejection-free for (C1)--(C4), and the resulting distribution is uniform over the
assignments satisfying those four constraints, because the weights are
path counts. When a live (C5) certificate is present the sampler is unchanged and
the completed assignment is tested against the literal disjunction
\eqref{eq:disj}, drawing again on failure, so the combined procedure is exact but
not rejection-free. The sampler
is used only off the deployed decision path, as the Monte-Carlo ground truth in
the E2 cross-checks of \S\ref{sec:results-verify}; the residual difference there
is at Monte-Carlo scale ($\numBeliefVsSampling$) rather than at the exact zero
E15 obtains by enumeration.

\subsection{Card groups and block groups}
\label{app:groups}

Every certificate \eqref{eq:disj} constrains only cards of one half-suit, so no
certificate couples two half-suits. That is what makes the exhaustive branch
affordable: enumerating one half-suit costs at most $\prod_{c \in H}|M_c|$
candidate assignments and is paid once per block, whereas handling the same $k$
live certificates by inclusion--exclusion over the whole of $U$ would cost $2^k$
forward--backward passes. The reference engine therefore partitions $U$
by half-suit and classifies each group:

\begin{description}[style=nextline]
  \item[Card group.] No live certificate in that half-suit. The group's cards are
        expanded one at a time by the recursions of \S\ref{app:layer}, with
        branching factor $|M_c| \le 5$. This is the ordinary capacity dynamic
        program restricted to the group.
  \item[Block group.] At least one live certificate. The group's assignments are
        enumerated exhaustively by depth-first search over the supports $M_c$ ---
        at most $5^6 = 15{,}625$ candidates for a six-card half-suit --- each is
        tested against every certificate of that half-suit, and the survivors are
        aggregated by count vector into the table $g_H(t)$ of
        \eqref{eq:blocktable}. Alongside $g_H(t)$ the engine accumulates, for
        each count vector, the per-card breakdown $g_H^{(c\to p)}(t)$ that
        \eqref{eq:q1} needs, so a single enumeration serves all three queries.
\end{description}

A group of either kind consumes a known number of cards, so the layer identity of
\S\ref{app:layer} survives the substitution and the outer dynamic program runs
over groups: $F_b(n) = \sum_t g_b(t) F_{b-1}(n+t)$, with $g_b(t)$ the indicator
of a single card placement in the card-group case. The path mass
\eqref{eq:pathmass} is then a contraction of the forward and backward arrays at
the group's entry layer, and the three queries \eqref{eq:q1}--\eqref{eq:q3} are
weighted sums over the group's own table. Nothing about this treatment is
approximate: the block enumeration is exhaustive, the certificate test is the
literal disjunction \eqref{eq:disj}, and the outer recursion is a counting
argument.

The exhaustive branch is what makes the equal-probability statement of
Corollary~\ref{cor:equalprob} usable and also what makes its caveat necessary.
The count-vector table records only survivors, so two allocations with the same
count vector are equiprobable \emph{given that both survive}; an allocation that
fails a certificate has probability zero even though its count vector may be
shared with survivors. The engine's \texttt{bestTeamAllocation} therefore returns
a stored survivor rather than reconstructing an allocation from the count vector.
E15 checks this directly: over \numOracleBestChecks{} calls, none returned an
allocation of probability zero and none returned a non-maximal one
(\S\ref{sec:results-verify}).

\subsection{The brute-force allocation oracle (E15)}
\label{app:oracle}

\begin{table}[!htbp]
\centering
\caption{Brute-force oracle (E15) against the exact reference engine.
Population: reachable states from \numOracleGames{} replayed \vfast{} versus
\vthree{} games, seed $20260822$, policy weights frozen; the unit is one state.
\numOracleStates{} states were enumerated, \numOracleStatesCFive{} with a live
ask-legality certificate, over \numOracleDeals{} consistent deals;
\numOracleStatesSkipped{} were skipped as too large. No intervals: every entry is
a deterministic maximum over the checks in the last column, except the sampler
row, a Monte-Carlo frequency comparison.}
\label{tab:oracle}
\small
\begin{tabular}{llrr}
\toprule
Quantity & Comparison & Value & Checks \\
\midrule
partition function $Z$ & Z max relative difference & $0$ & --- \\
per-card marginals $\mu_{c,p}$ & max absolute difference & $0$ & 656{,}826 \\
team ownership $\tau(H,T)$ & max absolute difference & $0$ & 26{,}930 \\
named allocation $\alpha(A)$ & max absolute difference & $0$ & 3{,}189{,}103 \\
equal-probability corollary & max within-class spread & $0$ & --- \\
sampler frequencies & max absolute difference & 0.0532 & 46{,}345{,}121 \\
\bottomrule
\end{tabular}
\end{table}


The block engine of \S\ref{sec:blocks} and the card-level dynamic program of
\S\ref{sec:counting} share a derivation, so agreement between them is weak
evidence that either enumerates the right collection of assignments. E15 supplies
independent evidence by recomputing the same posterior with no dynamic
programming, no factorisation and no sampling. The enumerator is
\filepath{engine/src/oracle.hpp} and the driver is the \texttt{oracle} command of
\filepath{engine/src/main.cpp}.

\paragraph{The enumerator.} For one observer's constraint state the unresolved
cards are ordered and searched by iterative depth-first search over their
assignments to seats. At depth $j$ the candidate seats for card $c_j$ are exactly
the surviving support $M_{c_j}$, which carries (C1)--(C3), and a candidate is
pruned at once if it would take the number of cards already given to that seat
above its capacity $q_p$, which carries (C4). Every leaf is therefore a complete
capacity-feasible assignment, and (C5) is tested there: the literal disjunction
\eqref{eq:disj} is evaluated against the completed assignment, certificate by
certificate, with no relaxation and no early aggregation. Surviving leaves are
counted into $Z$, tallied into per-card ownership counts, and tallied into a
per-half-suit table keyed by the packed assignment of that half-suit's unresolved
cards. Nothing is factorised: the count behind every named allocation is an
explicit count of leaves.

\paragraph{Coverage.} The search is exponential in $|U|$, so each state carries a
cut-off. A state is attempted only when the capacity dynamic program reports at
most \numOracleMaxDeals{} consistent deals, and the search is abandoned if the
product of the support sizes or the node count exceeds a fixed multiple of that
budget. A state that hits either bound is counted as skipped and reported ---
\numOracleStatesSkipped{} of them against \numOracleStates{} enumerated --- rather
than dropped silently, so the coverage of the test is visible rather than implied.
Skipping is correlated with the property that makes a state hard, a large
unresolved pool, so E15 is evidence about small reachable states and not about
the whole state space.

\paragraph{What is compared.} The block engine is built on the same constraint
state, and seven quantities are compared on every enumerated state: the partition
function $Z$; the per-card marginals \eqref{eq:q1} against the leaf tallies; the
team-ownership probabilities \eqref{eq:q3}, for both teams and every half-suit
still feasible for the team in question, against the summed counts of the
allocations lying wholly inside it; the probability \eqref{eq:q2} of \emph{every}
named allocation the enumeration produced rather than a sample of them; the
equal-probability corollary (Corollary~\ref{cor:equalprob}), as the spread
between the largest and the smallest enumerated probability within each
count-vector class; the reference declarer's own named allocation, by locating
\texttt{bestTeamAllocation}'s output in the enumerated table and asking both
whether it has positive probability and whether it attains the maximum over that
team's allocations; and the frequencies of the direct sampler of
\S\ref{app:sampler} drawn on the same state, against the enumerated marginals.
The command and its per-state draw budget are in \S\ref{app:repro-commands}.

\paragraph{Why the deterministic differences are exactly zero.} Six of the seven
comparisons are deterministic. Five of those are the zero entries of
Table~\ref{tab:oracle}, and the sixth, the reference declarer's allocation, is
reported in \S\ref{sec:results-verify} as \numOracleBestChecks{} checks with
\numOracleBestBad{} inconsistent. Zero here is not rounding luck. Both sides
compute integer counts over the same finite set of assignments --- the enumerator
by incrementing a counter at each surviving leaf, the block engine by the
counting recursions \eqref{eq:fwd}--\eqref{eq:Z} and the block tables
\eqref{eq:blocktable} --- and both hold those counts in double precision, where
the per-state cut-off of \numOracleMaxDeals{} consistent deals keeps every one of
them far below $2^{53}$ and therefore exactly represented. The two sides form the
same ratio
$\text{count}/Z$ from identical operands and obtain identical doubles, so the
difference is zero rather than small and the comparison has no tolerance
absorbing an error. The seventh comparison, the sampler, is a Monte-Carlo
frequency and differs at Monte-Carlo scale: \numOracleSampleDiff{} over
\numOracleDraws{} accepted draws.

\subsection{Declarations require no special handling}
\label{app:decl}

A declaration --- correct or incorrect --- removes six cards from play and
changes the public hand counts of the players who held them. Both effects are
public. In the constraint system this means the six cards leave $U$ with their
holders known, the affected half-suit leaves $\mathcal{H}$, and each $h_p$ in
\eqref{eq:capacity} drops by the number of those cards that seat held. The
capacities $q_p = h_p - k_p$ therefore re-derive correctly with no bespoke rule:
the same identity that absorbs a successful ask absorbs a declaration. The one
thing the observer learns beyond the removal is the declarer's belief, which the
rules posterior does not use; \S\ref{sec:policyprior} is where information of
that kind would enter.

The same argument explains why an incorrect declaration needs no special case
either. It reveals the true holders of all six cards, which is strictly more
information than a correct one gives the opposing team, and that information
arrives through exactly the same channel: resolved holders and updated public
counts.

\subsection{The Sinkhorn conditioning step}
\label{app:sinkhorn}

The Fast path maintains a matrix $p_{c,p}$ over unresolved cards and seats,
initialised to the prior weights \eqref{eq:prior} on the support $M_c$ and zero
off it. One Sinkhorn iteration normalises each row to sum to one, then scales
each column $p$ by $q_p / \sum_c p_{c,p}$, which enforces \eqref{eq:capacity} in
expectation. The deployed setting is four outer sweeps of eight such iterations,
with a conditioning step for the certificates applied between consecutive sweeps
and a final row normalisation.

The conditioning step derives as follows. Fix a certificate with asking seat $a$
(the notation of (C5); $A$ remains the named allocation of \S\ref{sec:notation})
and card set $D$, and let $E$ be the event $\bigvee_{d\in D}[\sigma(d)=a]$.
Treating the entries $\{p_{d,a}\}_{d \in D}$ as independent Bernoulli indicators,
\[
  \Pr[\neg E] \;=\; \prod_{e \in D} (1 - p_{e,a}),
  \qquad
  \Pr[E] \;=\; 1 - \prod_{e \in D} (1 - p_{e,a}) .
\]
For a particular $d \in D$, the event $\sigma(d) = a$ implies $E$, so
\[
  \Pr[\sigma(d) = a \mid E] \;=\; \frac{p_{d,a}}{\Pr[E]} .
\]
For a seat $p \ne a$, the event $\sigma(d) = p$ is compatible with $E$ only if
some other card of $D$ goes to $a$, whose independent probability is
$1 - \prod_{e \in D\setminus\{d\}}(1 - p_{e,a})$. Together the two cases give the
conditioning step of \S\ref{sec:fast},
\begin{equation}
  \Pr[\sigma(d) = a \mid E] = \frac{p_{d,a}}{1 - \prod_{e \in D}(1 - p_{e,a})},
  \qquad
  \Pr[\sigma(d) = p \mid E] = p_{d,p}\,
    \frac{1 - \prod_{e \in D \setminus \{d\}}(1 - p_{e,a})}{1 - \prod_{e \in D}(1 - p_{e,a})}.
  \label{eq:cond}
\end{equation}
Rows are renormalised afterwards, and
the next Sinkhorn sweep restores the capacity constraints that the reweighting
disturbed. The implementation guards two degenerate cases: when $\Pr[E]$
underflows, the mass is forced onto the certificate rather than dividing by a
near-zero denominator; when $\Pr[\neg E]$ underflows, the certificate is already
implied and the step is skipped.

The independence assumption is false --- the capacity constraints couple the
cards of $D$ --- so \eqref{eq:cond} is an approximation, and interleaving it with
capacity fitting does not converge to the exact marginals of
\eqref{eq:q1}. E2 measures the residual gap directly against the reference
engine: mean absolute difference \numSinkhornMean{}, maximum \numSinkhornMax{}
(\S\ref{sec:results-verify}).

\subsection{The deployed declaration estimator $\hat\alpha$}
\label{app:alpha}

The declaration query \eqref{eq:q2} asks for the joint probability of a named
allocation. The Fast path has only marginals, so it evaluates the joint by
sequential conditioning. Given a candidate allocation $A$ that assigns cards
$c_1,\dots,c_6$ of one half-suit to seats $p_1,\dots,p_6$, the estimator is
\[
  \hat\alpha(A) \;=\; \prod_{j=1}^{6} \hat\mu^{(j)}_{c_j, p_j},
\]
where $\hat\mu^{(j)}$ is the Fast marginal recomputed on the constraint state in
which $c_1,\dots,c_{j-1}$ have already been resolved to $p_1,\dots,p_{j-1}$. Each
step fixes one card, decrements that seat's capacity in \eqref{eq:capacity},
propagates the resulting support reductions, and refits. A card already resolved
contributes a factor of one if it matches $A$ and zero otherwise; a card whose
support excludes its assigned seat returns zero immediately; and the product
short-circuits once it falls below $10^{-9}$, since candidates that weak are
rejected regardless.

The chain rule this implements is exact --- the joint is the product of
successive conditionals in any fixed order --- so the estimator is exact in
structure and approximate only through the marginal solver of
\S\ref{app:sinkhorn}. Refitting between cards carries the capacity coupling of
\eqref{eq:capacity}, which a product of unconditioned marginals ignores; no
experiment in this paper measures the two estimators against each other, so the
comparison is structural rather than empirical.

Naming the allocation is a separate step, and the deployed default handles it
more crudely. For each card of the candidate half-suit, the seat is chosen as the
teammate with the largest \emph{unconditioned} Fast marginal $\hat\mu_{c,p}$,
with a card the observer holds assigned to the observer and a card already
resolved to a teammate assigned to that teammate. Only then is $\hat\alpha$
evaluated on the resulting allocation. Choosing greedily under unconditioned
marginals is not the same as maximising the joint, precisely because of the
coupling just described; the alternative of choosing each seat under the
\emph{conditioned} marginal at that step is implemented and is one of the
frozen-policy ablations of E5, where it changed the win rate by \numAblGmap{}
points with a 95\% paired percentile bootstrap interval over deals of
\numAblGmapCI{} (\S\ref{app:ablations}). Under
\vblock{} the naming step is replaced by \texttt{bestTeamAllocation}, which
returns a maximum-probability survivor directly from the block table
(\S\ref{app:groups}).

\subsection{Where each path is used}
\label{app:paths}

To summarise the division of labour. \vfast{} runs the Fast path of
\S\ref{sec:fast} for every query on the decision path: ask scoring, the two-ply
branch re-evaluation, the declaration pre-gates, and $\hat\alpha$. The reference
engine of \S\ref{sec:blocks} is invoked in three places only --- the cross-checks
of E2, the brute-force oracle comparison of E15, and the \vblock{} substitution
ablation --- none of which is part of the deployed policy. The throughput cost of
the difference is reported in \S\ref{sec:throughput}.
